%%%
\chapter{Expedient Trickle Sync Algorithm} 
%%%

\section{Introduction}  

The sole goal of the Expedient Trickle Sync (ETS) algorithm is to reduce the cost involved with data synchronization.  It does so by considering two factors:  the cost involved with physically transferring bits across a network, and the cost conceptually incurred when the user makes use of out-of-date data.  These two factors are frequently diametrically opposed:  the network cost can easily be reduced to zero by never synchronizing any information, however the cost involved with the use of out-of-date data increases exponentially.  Conversely, we can synchronize at a high rate of frequency to minimize the probability of using out-of-date information, however this can be at the expense of using more network time and resources than is otherwise necessary.  Thus, minimizing the cost of synchronization is a fine balance between ensuring the data the user is most likely to need is up-to-date, while minimizing the frequency and amount of time spent synchronizing.

The ETS algorithm attempts to minimize the cost of synchronization via four heuristic sub-algorithms:  determining when to synchronize, determining which available network to synchronize with, determining which subset of records to synchronize, and lastly by prioritizing the synchronization order.  While the intent of these four sub-algorithms is to reduce the overall cost of synchronization, they are not designed to guarantee absolute minimal synchronization cost; additionally, as they are based in part on a statistical analysis of the users past needs, sudden changes in a users habits could potentially circumvent these algorithms.  However, as these new user patterns take hold, the algorithm can learn from them in order to once again attempt to minimize the network cost; as such, the goal in cost reduction of the ETS protocol may not always be optimal when measured across short durations, it should approach an optimal level over the long term.
  
\section{Determining when to synchronize}

Determining when the synchronize so as to minimize cost is an extremely complex task.  At a one-second resolution, there are $2^{86400}$ possible sets of synchronization points.  Determining which members of these sets have the minimal cost in a rigourous manner would be inordinately computationally intensive, and thus is completely infeasible.  As such, a heuristic mechanism which will quickly reduce the cost of synchronization with minimal computation is desirable.  Such a heuristic must be suitable for running on a mobile device, and thus cannot require significant computing power or memory.

To fulfill these requirements, an algorithm which we can very quickly evaluate is important.  Easily calculated and recorded values, such as those pertaining to the probability the user will need a record on the device within a specified interval of time, how long it has been since the device was last synchronized, and the cost of synchronizing on the least expensive network in the mobile devices current location are excellent candidates for creating such an algorithm.

The mechanism selected for the ETS algorithm to use is to combine the above mentioned values into an equation, such that high probabilities and time since last synchronization have an increasing influence on the expression, and where high cost has a decreasing influence on the expression.  The derived value is then compared against a threshold value ($T$):  if the value is greater than the threshold at the time it is computed, we initiate a synchronization session, otherwise we do not.

In specific, given the probability function $P(x)$ which generates the probability that the user will access a record on the mobile device during the fixed-size interval of the day $x$, $\Delta t$ is the time since the last synchronization session, and $C_{byte_{net}}$ is the cost of transferring one byte of data on the selected network, then we define the function $G$ as:
\begin{equation}
G(x) = \frac{(\frac{8}{15}*P(x+1)+\frac{4}{15}*P(x+2)+\frac{2}{15}*P(x+3)+\frac{1}{15}*P(x+4)) * \Delta t}{C_{byte_{net}}}
\end{equation}\label{threshold_function}
Note that in the function $G$ we evaluate the mobile device record access probabilities of the next four time intervals, with each subsequent interval having half the influence of its immediately preceding interval.  In the default implementation, the intervals are 15 minutes in length; thus the probability calculation takes into consideration the next hours worth of usage.

Calculation of these probabilities is based upon a table maintained by the ETS algorithm, which stores the number of record accesses for each time interval in the day, for each of the last seven days.  In order to keep computation simple, the quantity of records accessed within each interval is not considered, only whether or not \emph{any} record was accessed during the interval; the probability algorithm is interested only in the probability of a record being accessed each day during the specified interval, and not how many were accessed during each interval.  Thus, an interval $x_a$ which has accesses in the seven preceding days of $\{1, 1, 1, 1, 1, 1, 1\}$ would have $P(x_a) = 1.0$, whereas an interval $x_b$ which has access counts in the seven preceding days $\{7, 0, 0, 0, 0, 0, 0\}$ would have $P(x_b) =  \frac{1}{7}$.  Using only the last seven days of access information permits the system to adapt within a one week interval to changing record access habits by the user.

Given function \ref{threshold_function}, and given a function $I$ which returns the fixed-sized time interval containing a specific instant of time $t$, we can determine whether or not to synchronize by evaluating the following inequality:
\begin{equation}
G(I(t)) > T
\end{equation}\label{threshold_inequality}
This inequality is evaluated every second during the day to determine whether or not synchronization needs to occur.  Given a sufficiently low probability, a sufficiently low time since the last synchronization, a sufficiently high network transfer cost, or any combination of these three attributes, synchronization will not occur during the interval.  Given a sufficiently high probability, long time since the last synchronization, or low network transfer cost, or any combination of the three, synchronization will occur.

It should be noted that synchronization when the network transfer cost is low presents certain potential problems.  As the data transfer cost for a given network approaches zero, synchronization will always occur, and the time since the last synchronization ($\Delta t$) would always be equal to 1 second (as we would synchronize during every interval due to the minor cost impact synchronization would have in a "free network" environment).  ETS prevents a divide-by-zero situation simply by avoiding it altogether -- even a "free" network is assumed to have some minor cost for the electrical power to perform the transfer, however for a mobile device synchronizing during every 1s interval while in a "free" network zone is most likely impossible, and would be a significant battery drain.  As such, the ETS algorithm prevents this situation from occurring by only permitting one synchronization during each fixed sized interval $x$ so as to prevent an unreasonable use of resources by the data synchronization subsystem.

\subsection{Evaluating the Threshold}

The algorithm as described is heavily reliant upon the value of the threshold to influence the frequency of synchronization of the mobile devices data, and as such selection of the threshold value will have a direct impact on the cost of synchronization.  As usage patterns and networks differ from user to user, it is impossible to pre-calculate the threshold value which specifically reduces cost; as such, a heuristic mechanism is used to "learn" an optimal value for the threshold which reduces overall cost of synchronization.



*** More stuff goes here

\section{Determining which network to synchronize via}

Determining which network to use for synchronization is the simplest sub-algorithm involved in the ETS algorithm.  When the mobile device is in the presence of multiple networks, and the decision has been made to synchronize, the network with the lowest cost per byte factor is selected for use during the synchronization session.  If multiple networks with the same cost per byte are available, then the network with the fastest transfer rate at that cost is selected.  If multiple networks with the same cost and transfer rate are present, one is selected at random.  Real-world implementations of the ETS protocol may decide to use further criteria for network selection (such as prioritizing those on a secure or trusted network).

\section{Determining which records to synchronize}

Once a synchronization session has been initiated, the ETS protocol must decide which records to synchronize.  Unlike most existing data synchronization algorithms, the ETS algorithm doesn't always attempt to synchronize every modified record during the synchronization session, but instead uses a  variety of metrics to balance the real cost of synchronizing each record versus the expected intangible cost of not having said record up-to-date.

Given the set of records needing synchronization $R$, with size $|R| = n$ and individual records $r_i \in R$, where $i$ is the records index value in the set of records to be synchronized, the function $C$ that calculates the cost of transferring a record via the selected network, the function $P$ which calculates the probability that a record will be accessed by the user, the function $V$ which calculates the version number of a given record on a given device (where $s$ is the server and $h$ is the mobile device), the constant $k$ denoting the expected cost of using out-of-date information, and decision variables $x_i \in \{0, 1\}$, the expected cost $XC$ of the synchronization can be expressed as:
\begin{equation}
XC = \sum_{i=1}^n C(r_i)x_i+P(r_i)\overline{x_i}\cdot(V(r_{i_s})-V(r_{i_h}))\cdot k
\end{equation}\label{expected_sync_cost}

Thus the goal of the record selection sub-algorithm is to minimize $XP$ by solving for all $x_i$.  This can be done by testing the following equality for each $r_i$ requiring synchronization, and evaluating the corresponding $x_i = 1$ when it holds true:
\begin{equation}
C(r_i) < P(r_i)\cdot(V(r_{i_s})-V(r_{i_h}))\cdot k
\end{equation}\label{sync_cost_inequality}

In essence, we are directly comparing the cost of synchronizing the records, versus the expected cost of not synchronizing the record.  If the real network cost of synchronizing the record is higher than the cost of not having the record available, the ETS algorithm will not synchronize it during this session.  If the real cost of synchronizing the record is lower than the cost of not having the record available, the record will be synchronized.

In practice, the above inequality implies that when we are using a sufficiently expensive network, not all records will necessarily be synchronized, but instead only those records which either have a high expectation of use by the user, which are sufficiently out-of-date, or a combination of these two factors.  When using a sufficiently inexpensive network however, virtually all records will be expected to be synchronized.

The key to the proper functioning of this sub-algorithm is the suitable selection of a value for $k$.  If $k$ is set too high, the sub-algorithm will synchronize every record during every session, whereas if $k$ is set too low, it's possible that records will never be synchronized.  As the value of $k$ is application and domain specific, it is impossible for the ETS algorithm to determine this value on its own, and thus must be obtained from an external source, such as the user or an administrator.

\section{Prioritizing synchronization order}

Finally, as its last step prior to exchanging records between the server and the mobile device, the ETS algorithm re-orders the records to be synchronized based on their probability $P$ of access.  This is done in an attempt to ensure that, should a network disconnection occur at synchronization time $t$, the records yet to be transferred will have a lower access probability than the records which have already been transferred (if any).  This is performed as an aid in cost reduction when network disconnections are proportionately high; there is no cost savings involved with this sub-algorithm when all data that has been selected for transfer during the current synchronization session is successfully exchanged.

\section{Comments}

*** Add final comments on the ETS algorithm here ***

\subsection{Problems with Thresholding}

*** Talk about problematic patterns here, and how we deal with them

